% TOP_PROBLEM: {"id": 3273, "title": "Large acyclic induced subdigraph in a planar oriented graph.", "category_id": 3, "category": {"id": 3, "name": "graph_theory", "display_name": "Graph Theory", "description": "Problems involving graphs, networks, and their properties.", "slug": "graph-theory", "order_index": 3, "created_at": "2026-07-31T15:26:25.671Z"}, "status": "open", "problem_number": "OPG-52197", "set_id": 12}
% FIRST_POSTED: 2026-10-01
% ENTRY_KIND: statement_only
% UnsolvedMath ID 3273; source catalogue code OPG-52197
% !TeX program = lualatex
% Definition and sources only; this file is not an LLM solution attempt.
\documentclass[11pt,a4paper]{article}
\usepackage[margin=25mm]{geometry}
\usepackage{fontspec}
\setmainfont{lmroman10-regular.otf}[BoldFont=lmroman10-bold.otf,
 ItalicFont=lmroman10-italic.otf,BoldItalicFont=lmroman10-bolditalic.otf]
\usepackage{amsmath,amssymb,mathtools}
\usepackage{unicode-math}
\setmathfont{latinmodern-math.otf}
\AtBeginDocument{\let\setminus\smallsetminus}
\usepackage{xurl,hyperref}
\hypersetup{colorlinks=true,urlcolor=blue}
\setcounter{secnumdepth}{0}
\setlength{\parindent}{0pt}
\setlength{\parskip}{5pt}
\setlength{\emergencystretch}{3em}
\begin{document}
\section*{Large acyclic induced subdigraph in a planar oriented graph.}\label{3273}
{\small UnsolvedMath ID 3273 · OPG-52197 · Graph Theory · OpenGarden}
\subsection{Definitions and mathematical statement}
Let $D=(V,A)$ be a finite oriented graph whose underlying simple
undirected graph is planar. Thus there are no loops, repeated arcs,
or opposite pairs of arcs. Put $n=|V|$. For $U\subseteq V$, the induced
subdigraph $D[U]$ contains all arcs of $D$ having both ends in $U$.
It is acyclic when it has no directed cycle, equivalently when its
vertices admit a linear ordering in which every arc points forward.

The conjecture asks whether one can always find $U\subseteq V$ with
\[
                           |U|\ge\frac35n
                  \quad\text{and}\quad D[U]\text{ acyclic}.
\]
Since $|U|$ is an integer, the size condition means
$|U|\ge\lceil3n/5\rceil$. The equivalent deletion formulation asks for
a set $S\subseteq V$ of at most $\lfloor2n/5\rfloor$ vertices meeting
every directed cycle; then $U=V\setminus S$. Such an $S$ is a directed
feedback vertex set.

Deleting arcs is not allowed when testing whether $D[U]$ is acyclic:
all arcs between retained vertices remain. On the other hand,
undirected cycles are allowed if their orientations do not form
directed cycles. This distinguishes the target from an induced-forest
problem for the underlying graph. It also asks for one large acyclic
set, not a partition of all vertices into a fixed number of acyclic
classes. The graph may be disconnected, and the orientation is arbitrary
subject to planarity of the underlying graph.
\subsection{Short English statement}
Can at least three fifths of the vertices of every planar oriented graph be retained without retaining any directed cycle?
\subsection{Sources}
\begin{itemize}
\item[S1] ULAM.AI and the UnsolvedMath Contributors, supplied problems.json, record 3273 (OPG-52197), OpenGarden collection. Catalogue identity and source transcription; CC BY 4.0.
\url{https://huggingface.co/datasets/ulamai/UnsolvedMath}.
\item[S2] Open Problem Garden, Large acyclic induced subdigraph in a planar oriented graph.. Original problem and discussion; the mathematical formulation below is expanded from the supplied source record.
\url{http://www.openproblemgarden.org/op/large_acyclic_induced_subdigraph_in_a_planar_oriented_graph}.
\end{itemize}
\end{document}
